\section*{Capítulo \ref{ch:b3:covariant-electromagnetism} --- Eletromagnetismo covariante}

\begin{solution}{exo:b3:covariant-electromagnetism:1}
(a) $\gamma = 1.25$: $a'^0 = 1.25(5 - 1.8) = 4.0$, $a'^1 = 1.25(3 -
3) = 0$; $a\cdot a = 25 - 9 = 16 = 16 - 0$. (b) A transformação é
linear e, portanto, se distribui sobre somas. (c) $(c, \vect v)$: não
--- sua componente ``temporal'' é a mesma para todas as partículas,
enquanto a mistura exige o contrário; $\gamma(c, \vect v) = \dd x^\mu/\dd\tau$:
sim, um \omterm{def:b3:covariant-electromagnetism:fourvector}{quadrivetor} dividido pelo invariante $\dd\tau$. (d) As três
componentes de $\vect E$ se misturam com as de $\vect B$, e não com
escalar algum: os campos preenchem um tensor de dois índices, não um
\omterm{def:b3:covariant-electromagnetism:fourvector}{quadrivetor}.
\end{solution}

\begin{solution}{exo:b3:covariant-electromagnetism:2}
(a) $u = I/nSe = 10/(\num{8.5e28} \times \num{e-6} \times
\num{1.6e-19}) = \qty{7.4e-4}{m/s}$. (b) Rede: $(nec, \vect 0)$;
elétrons: $(-nec, -ne\vect u)$ --- a densidade de corrente deles,
$\rho_-\vect u = -ne\vect u$, aponta contra a sua deriva, e essa
\emph{é} a direção da corrente convencional. (c) Ambas são estáticas e uniformes: todos os termos se anulam. (d)
Total: $(0, \vect\jmath)$ com $j = I/S = \qty{e7}{A/m^2}$ --- uma
corrente pura sem carga, a configuração que torna sutil a relatividade
do fio.
\end{solution}

\begin{solution}{exo:b3:covariant-electromagnetism:3}
(a) $E_y$: $F^{20} = E_y/c$; $B_x$: $F^{32} = B_x$. (b) Para $\vect B
= B\vect e_z$, apenas $F^{12} = -B$ e $F^{21} = +B$; para a onda,
$F^{20} = E/c$, $F^{12} = -E/c$ (e os parceiros antissimétricos). (c)
$F^{12} = \partial^1A^2 - \partial^2A^1 = -\partial_xA_y +
\partial_yA_x = -(\operatorname{\vect{curl}}\vect A)_z = -B_z$:
confere com a matriz. (d) A taxa de variação da energia da partícula é
$\propto F^{0\nu}u_\nu$; no referencial de repouso $u_\nu = (c, \vect 0)$
e a antissimetria faz $F^{00} = 0$: uma força que jamais pode realizar
trabalho sobre uma partícula em repouso --- a propriedade que define a
parte magnética.
\end{solution}

\begin{solution}{exo:b3:covariant-electromagnetism:4}
(a) Impulso ao longo do campo: componentes longitudinais intactas, $\vect
E' = E_0\vect e_y$, $\vect B' = \vect 0$. (b) Impulso ao longo de $\vect
e_x$: $E'_y = \gamma E_0$, $B'_z = -\gamma vE_0/c^2$ --- as placas
passam a escoar como correntes superficiais $\pm\sigma'v$, e duas
lâminas de corrente opostas envolvem exatamente um campo desses. (c) $E'^2 - c^2B'^2 =
\gamma^2E_0^2(1 - \beta^2) = E_0^2$; $\vect E'\cdot\vect B' = 0$: os
dois se preservam. (d) As placas se contraem ao longo de $x$, logo $\sigma' =
\gamma\sigma$, coerente com $E' = \sigma'/\varepsilon_0 = \gamma
E_0$; a separação, transversal ao impulso, fica intacta.
\end{solution}

\begin{solution}{exo:b3:covariant-electromagnetism:5}
(a) A força magnética $-e\vect v\wedge\vect B$ empurra os elétrons ao
longo da haste: fem $= vBL$. (b) No referencial da haste, ela está em
repouso --- nenhuma força magnética sobre cargas paradas --- mas a
transformação entrega $\vect E' = \gamma\,\vect v\wedge\vect B$: é um
campo elétrico honesto que faz o trabalho. (c) $E'L = \gamma vBL
\approx vBL$ em ordem $v/c$. (d) Tinha de funcionar porque a fem ``de
movimento'' e a fem ``de transformador'' são um só fenômeno lido em
dois referenciais: a regra do fluxo é a covariância vestida de 1831 ---
o artigo de Einstein de 1905 abre exatamente com essa assimetria entre
o ímã e o condutor.
\end{solution}

\begin{solution}{exo:b3:covariant-electromagnetism:6}
(a) As componentes $x$ ficam intactas; para as transversais,
$E_y'^2 + E_z'^2 - c^2(B_y'^2 + B_z'^2) = \gamma^2[(E_y - vB_z)^2 +
(E_z + vB_y)^2 - c^2(B_y + vE_z/c^2)^2 - c^2(B_z - vE_y/c^2)^2]$;
os termos cruzados se cancelam e os quadrados reúnem $\gamma^2(1 -
\beta^2) = 1$ vezes a combinação não transformada. (b) A mesma
contabilidade sobre $E_xB_x + E_yB_y + E_zB_z$. (c) $E^2 - c^2B^2 =
3c^2B^2 > 0$: impulso ao longo de $\vect E\wedge\vect B$ a $v = c^2B/E =
c/2$; ali $B' = 0$ e o campo que sobra é $E' = E/\gamma =
\sqrt3\,cB$ --- cujo quadrado dá o invariante, como deve ser. (d) Uma onda de luz tem os dois invariantes
nulos: todo referencial precisa reproduzir $E' = cB' \perp$, jamais um campo puro.
\end{solution}

\begin{solution}{exo:b3:covariant-electromagnetism:7}
(a) Com $v_{\text{d}} = E/B$: $E'_y = \gamma(E - v_{\text{d}}B) =
0$; $B'_z = \gamma(B - v_{\text{d}}E/c^2) = B/\gamma$: campo
puramente magnético, um pouco enfraquecido. (b) Ali: um círculo a
$qB'/\gamma m$; de volta ao laboratório: esse círculo mais a deriva
uniforme --- uma cicloide que rasteja a $E/B$ perpendicularmente aos
dois campos, a trajetória das cargas num magnétron. (c) Uma partícula
que se move exatamente a $v_{\text{d}}$ está em repouso no referencial
de deriva, onde o único campo é magnético e ela nada sente: não é
desviada. (d) Para $E > cB$ a deriva necessária excederia $c$: não
existe tal referencial; em vez disso existe um referencial de $E$ puro
(exercício anterior), e a carga é acelerada sem limite ao longo do
campo.
\end{solution}

\begin{solution}{exo:b3:covariant-electromagnetism:8}
(a) A fase $\omega t - \vect k\cdot\vect r$ conta eventos de passagem
de crista, um número invariante; ela é igual a $k^\mu x_\mu$, e a
invariância do produto para todo $x^\mu$ obriga $k^\mu$ a se
transformar como \omterm{def:b3:covariant-electromagnetism:fourvector}{quadrivetor}. (b) $\omega'/c = \gamma(\omega/c -
\beta k_x)$ com $k_x = (\omega/c)\cos\theta$; para $\theta = 0$,
$\omega' = \omega\sqrt{(1-\beta)/(1+\beta)}$. (c) $\cos\theta' =
k'_x/(\omega'/c)$: a fórmula enunciada. (d) $\beta = \num{e-4}$:
$20.5''$; ao longo de um ano a posição aparente varre uma elipse desse
semieixo angular --- a aberração de Bradley.
\end{solution}

\begin{solution}{exo:b3:covariant-electromagnetism:9}
(a) No referencial de repouso, Coulomb; o impulso multiplica as
componentes transversais por $\gamma$ (no plano transversal,
$E'_\perp = \gamma
q/4\pi\varepsilon_0b^2$), ao passo que o campo ao longo do movimento,
avaliado à frente ou atrás a uma distância $r$ do laboratório,
corresponde a uma distância $\gamma r$ no referencial de repouso:
reduzido por $1/\gamma^2$. (b) O disco de largura angular
$1/\gamma$ passa com velocidade $v$: duração $\sim (b/\gamma)/v$. (c)
$E' = 7250 \times \num{1.6e-19} \times \num{9e9}/\num{e-4} \approx
\qty{0.1}{V/m}$, durando $2b/\gamma c \approx \qty{9}{fs}$. (d)
$B' = vE'/c^2 \approx E'/c$ e os dois invariantes são $O(1/\gamma^2)$:
localmente indistinguível de um pulso de luz --- a base da descrição
por ``fótons equivalentes'' das colisões de cargas rápidas.
\end{solution}

\begin{solution}{exo:b3:covariant-electromagnetism:10}
(a) $E'_y = \gamma(E - vB) = \gamma E_0(1 - \beta)\cos(\cdots) =
E_0\sqrt{(1-\beta)/(1+\beta)}\cos(\cdots)$, e $B'_z = E'_y/c$ pela
mesma álgebra. (b) A fase se transforma pelo mesmo fator Doppler: a
mesma onda nula, mais vermelha e mais fraca. (c) A densidade de energia
cai com o quadrado do fator Doppler; o número de fótons não muda ---
cada $h\nu$ de fóton carrega o fator inteiro. (d) ``Viajar ao lado''
significa o fator $\to 0$ quando $\beta \to 1$: a onda nunca se torna
estática porque seus invariantes são nulos --- não há referencial em
que a luz fique parada, que é onde a relatividade deste livro começou.
\end{solution}

\begin{solution}{exo:b3:covariant-electromagnetism:11}
(a) Por volta, $\Delta E = P \times 2\pi r/c$: o resultado enunciado.
(b) $\gamma = \num{e11}/\num{5.11e5} = \num{1.96e5}$: $\Delta E \approx
\qty{2.9}{GeV}$ por volta --- três por cento da energia do feixe
reinjetados a cada volta; os clístrons do LEP eram o maior transmissor
de rádio da Terra. (c) Prótons: $\gamma = 7250$, com $\gamma^4$ menor
por $(m_{\text{p}}/m_{\text{e}})^4 \approx \num{e13}$: cerca de
\qty{5}{keV} por volta --- desprezível. (d) $\gamma^4 = (E/mc^2)^4$: a
energias iguais o elétron irradia $\num{e13}$ vezes mais; daí os
colisores lineares (irradia-se uma vez, não a cada volta) ou anéis
gigantescos para elétrons, ao passo que os anéis de prótons escalam sem
dor até \qty{27}{km}.
\end{solution}

\begin{solution}{exo:b3:covariant-electromagnetism:12}
(a) $|\vect p|$ é constante (não há trabalho); $\dot{\vect p} = q\vect
v\wedge\vect B$ gira $\vect p$ à taxa $qvB/p = qB/\gamma m$. (b)
Escorregamento por volta $2\pi(\gamma - 1)$; com $\gamma - 1$ crescendo
linearmente até $\Gamma$, o escorregamento total é
$\approx \pi N\Gamma$; um quarto de período é $\pi/2$: $\Gamma \approx 1/2N$.
Para $N = 100$: $\Gamma =
\num{5e-3}$, $E_k \approx \qty{4.7}{MeV}$ --- a escala certa; as
máquinas reais o esticaram até $\sim\qty{20}{MeV}$ com tensões altas
nos dês (menos voltas). (c) Sincrocíclotron: baixe a radiofrequência
durante cada pulso para acompanhar $qB/\gamma m$. Cíclotron isócrono:
deixe $B(r)$ crescer $\propto\gamma(r)$, de modo que a razão nunca mude
e o feixe permaneça contínuo. (d) $\gamma = 1 + 590/938 = 1.63$: o
campo na borda tem de superar o campo central por esse fator.
\end{solution}

\begin{solution}{pb:b3:covariant-electromagnetism:1}
\textbf{1.} $u = I/nSe = \qty{7.4e-4}{m/s}$; $u^2/c^2 =
\num{6e-24}$.
\textbf{2.} $\lambda = I/u = nSe = \qty{1.36e4}{C/m}$ --- catorze
quilocoulombs por metro em cada corrente.
\textbf{3.} $B = \mu_0I/2\pi d = \qty{2.0e-4}{T}$; $F = qvB =
\num{e-6} \times \num{e5} \times \num{2e-4} = \qty{2e-5}{N}$,
dirigida \emph{para} o fio (correntes paralelas se atraem).
\textbf{4.} As duas cargas lineares se cancelam exatamente:
$\lambda_{\text{tot}}
= 0$, nenhum campo em ordem zero.
\textbf{5.} Um elétron em excesso por metro: $E = 2k\lambda_1/d =
\qty{2.9e-7}{V/m}$, força $\qty{2.9e-13}{N}$ --- $10^8$ vezes menor
que a força magnética. O fio poderia esconder cem milhões de elétrons
avulsos por metro antes que a eletrostática rivalizasse com o
magnetismo: atribuir a força a $\vect B$ é seguro.
\textbf{6.} $\qty{2e-5}{N}$ são vinte pesos de grão de areia:
facilmente mensurável.
\textbf{7.} Rede: $-v$. Elétrons (que derivam a $u$ em sentido oposto
a $I$, isto é, oposto ao impulso): velocidade $(u + v)/(1 + uv/c^2)$
--- mais rápidos que a rede.
\textbf{8.} $\lambda' = \gamma_v(\lambda_{\text{tot}} - vI/c^2) =
-\gamma_v vI/c^2$.
\textbf{9.} $\lambda' = -\num{e5} \times 10/\num{9e16} =
\qty{-1.1e-11}{C/m}$: cerca de $7 \times 10^7$ elétrons em excesso por
metro. A corrente de elétrons, mais rápida neste referencial, é a mais
densa --- cada corrente se contrai pelo seu próprio $\gamma$.
\textbf{10.} $E' = 2k|\lambda'|/d = 2 \times \num{9e9} \times
\num{1.1e-11}/0.01 = \qty{20}{V/m}$, apontando para o fio; força
$qE' = \qty{2e-5}{N}$, atrativa.
\textbf{11.} Idêntica a $qvB$ a menos do fator $\gamma_v = 1 +
\num{5.6e-8}$: exatamente a lei de transformação de uma força
transversal (a quadriforça), $F_{\text{rep}} = \gamma F_{\text{lab}}$.
\textbf{12.} O fio continua conduzindo uma corrente (maior), e por
isso $\vect
B' \neq 0$ --- mas nossa carga está em repouso aqui, e um campo
magnético só agarra cargas em movimento.
\textbf{13.} $|\lambda'|/\lambda = \gamma_v vu/c^2 = \num{8.2e-16}$.
\textbf{14.} Ele multiplica $\lambda/e \approx \num{8.5e22}$ cargas
elementares por metro: $\num{8.2e-16} \times \num{8.5e22}
\approx \num{7e7}$ elétrons por metro --- macroscópico. A matéria é
tão enormemente carregada que sua neutralidade é o fio de uma navalha;
a relatividade inclina a navalha.
\textbf{15.} $F/L = \mu_0I_1I_2/2\pi d = \num{2e-7} \times 100/0.01 =
\qty{2e-3}{N/m}$ --- a fórmula que um dia \emph{definiu} o ampère.
\textbf{16.} $B = \mu_0nI = 4\pi\times10^{-7} \times \num{e4} \times
10 = \qty{0.13}{T}$; pressão magnética $B^2/2\mu_0 \approx
\qty{6300}{Pa} \approx \qty{0.6}{N/cm^2}$: carros inteiros pendurados
em relatividade integrada.
\textbf{17.} $c \to \infty$ mata $\lambda'$ e, com ele, a força, a
parte magnética do campo, os motores, os dínamos e os guindastes de
ferro-velho: o magnetismo não tem existência não relativística.
\textbf{18.} Em repouso, a carga vê o fio neutro do laboratório: as
duas correntes se cancelam com a precisão com que a matéria é neutra
--- conhecida experimentalmente com exatidão fantástica --- e força
alguma age.
\textbf{19.} Que a imagem das ``fontes dadas'' era a fatia, num único
referencial, de um só tensor: $\vect B$ é a parte do campo
eletromagnético que o movimento de um dado observador atribui a
correntes, e não a cargas.
\textbf{20.} O spin do elétron: cada elétron é um ímã elementar (uma
``corrente'' intrínseca e quântica), e o ferro é a matéria em que eles
se alinham --- o \cref{ch:b3:spin-two-level} e o
\cref{ch:b3:electromagnetism-in-matter} retomam isso.
\textbf{21.} No referencial de repouso do feixe a repulsão é
puramente coulombiana e \emph{máxima}; no laboratório, a atração
magnética das correntes paralelas quase a cancela. Sem contradição: a
força transversal do laboratório é a do referencial de repouso dividida
por $\gamma$, e a dinâmica mais lenta (\omterm{prop:b3:relativistic-kinematics:dilation}{dilatação do tempo}) completa a
contabilidade.
\textbf{22.} Força líquida no laboratório $= qE_{\text{lab}} - qvB_{\text{lab}} =
(1 - \beta^2)qE_{\text{lab}} = F_{\text{Coulomb}}/\gamma^2$, conforme a
transformação; no referencial de repouso o fator é transparente:
eletrostática pura, com a expansão observada através de um tempo
dilatado.
\textbf{23.} O vetor de Poynting: a energia escoa pelos \emph{campos}
em torno dos condutores a quase $c$, com os elétrons apenas a
ordenando; o cobre do fio é um guia, não um cano.
\textbf{24.} $vu/c^2 \approx \num{8e-21}$ --- e, ainda assim, Ørsted
viu sua bússola girar ao lado de um fio em 1820: o multiplicador de
$10^{22}$ cargas já estava em ação um século antes que alguém pudesse
nomeá-lo.
\textbf{25.} Um fio neutro de \qty{10}{A}, visto de \qty{e5}{m/s},
carrega $\qty{-1.1e-11}{C/m}$: a diferença de contabilidade de
$10^{-16}$ devida à \omterm{prop:b3:relativistic-kinematics:contraction}{contração dos comprimentos}, multiplicada por
$10^{23}$ cargas por metro, \emph{é} a força magnética --- todo o
magnetismo é a relatividade auditada a passo de caracol.
\end{solution}
